\documentclass[11pt,a4paper]{article}

\usepackage[T1]{fontenc}
\usepackage[utf8]{inputenc}
\usepackage{lmodern}
\usepackage{amsmath,amssymb,amsthm}
\usepackage{mathtools}
\usepackage{hyperref}
\usepackage{doi}
\usepackage{geometry}
\usepackage{booktabs}
\usepackage{array}
\usepackage{enumitem}

\input{glyphtounicode}
\pdfgentounicode=1

\geometry{margin=2.5cm}

\hypersetup{
  colorlinks=true,
  linkcolor=black,
  citecolor=black,
  urlcolor=black,
  pdftitle={The Spectral Gravity Sub-Programme (Presentation Note 4)},
  pdfauthor={Jérôme Beau},
  pdfsubject={Presentation note of the spectral gravity sub-programme of the Cosmochrony corpus},
  pdfkeywords={spectral gravity, heat-kernel expansion, Einstein tensor, spectral entropy,
    renormalization matching, Born-Infeld completion, graviton dispersion, capacity axiom}
}

\theoremstyle{definition}
\newtheorem{remark}{Remark}

\newcommand{\SPI}{\mathcal{S}_\Pi}
\newcommand{\Ag}{A_g}
\newcommand{\lsp}{\ell_{\mathrm{sp}}}
\newcommand{\GN}{G_{\mathrm{N}}}
\newcommand{\cEH}{c_{\mathrm{EH}}}
\newcommand{\cBI}{c_{\mathrm{BI}}}

\title{The Spectral Gravity Sub-Programme\\[0.4em]
\large Presentation Note 4}

\author{J\'{e}r\^{o}me Beau\\[0.3em]
\small Independent Researcher\\
\small \texttt{jerome.beau@cosmochrony.org}}

\date{2026}

\begin{document}

  \maketitle

  \begin{abstract}
    The spectral gravity sub-programme of the Cosmochrony corpus addresses a single
    central question: under which conditions does the projective spectral entropy functional
    $\SPI[g] = \tfrac{1}{2}\log\det' A_g$ support an infrared Einstein sector for a supplied
    Riemannian metric?
    The sub-programme separates the physical proper-time cutoff (power-sensitive terms
    $a_0\Lambda^4$, $a_2\Lambda^2$) from zeta regularization (finite and logarithmic terms
    governed by $a_4$), decomposes the renormalized metric variation into Einstein,
    cosmological, higher-derivative, and non-local sectors, and states the conditions
    under which the Einstein term dominates.
    The renormalized couplings --- including the sign and value of Newton's constant --- are
    matching data, not predictions.
    Conditionally on supplied coefficients, and at leading local infrared order, stationarity of the difference
    functional $\Gamma_\Pi = \SPI^{\mathrm{ren}} - W_\Pi$ yields the Einstein--matter coupling
    $G_{\mu\nu} + \Lambda_{\mathrm{ren}} g_{\mu\nu} = 8\pi\GN T_{\mu\nu}$.
    On the posited local four-derivative Lorentzian truncation, the massless graviton branch is
    exactly luminal, and the ultraviolet completion of determinantal Born--Infeld type is an
    admissible member of an explicitly parametrised family, not a uniquely selected one.
    The metric is supplied: its reading as the emergent metric of the programme is conditional on the
    unestablished spatial-limit hypothesis [H-L], the open lifting hypothesis [H-lift], a supplied rank-four
    extension and a continuation from Riemannian signature that is not performed, and no metric coefficient value is
    established.
    The sub-programme is constituted by four papers: Gravity, Thermodynamics,
    Lorentz/CausalPropagation, and BornInfeld.
    \textbf{Interpretive status.}
    The sub-programme establishes which local structures spectral geometry allows and where
    every remaining gap lies --- the matching principle for the couplings, the integrability
    bridge to any thermodynamic reading, the in-in construction of a Lorentzian spectral
    action, and the selection principle for the ultraviolet completion.
    Gravity emerging as a spectral response remains a structural possibility, not an
    established result.
  \end{abstract}

  \medskip
  \noindent\textbf{Keywords.}
  spectral gravity; heat-kernel expansion; Seeley--DeWitt coefficients; Einstein tensor;
  spectral entropy; renormalization matching; Born--Infeld completion; graviton dispersion;
  capacity axiom; presentation note.

  \newpage
  \tableofcontents

  \newpage

  \section{Motivation and Central Question}
    \label{sec:motivation}

    The spectral gravity sub-programme addresses the following question.

    \medskip
    \noindent\textit{Central question.}
    The emergent geometry sub-programme (Note~2~\cite{Beau2026egnote}) addresses the reconstruction of an
    effective metric from the admissibility structure.
    Under the unestablished spatial-limit hypothesis [H-L] and the open lifting hypothesis [H-lift] it extracts a
    rank-three
    co-metric (Q5b~\cite{Beau2026q5b}).
    Its target, a co-metric proportional to $\eta^{\mu\nu}$, holds only for a supplied rank-four extension
    satisfying the hypotheses of Q11~\cite{Beau2026q11} Corollary~6.1 (the conditions [T] and [INV$^+$] of
    Q8~\cite{BeauQ8},
    invariance of the four-dimensional principal symbol under the supplied cotangent action [ACT$_\tau$], and a
    negative ordering term), after rescaling $\tau$, with
    its overall scale free; no metric coefficient value is established.
    The question of the dynamics is independent of that reconstruction, and the present sub-programme poses it for
    a supplied Riemannian metric $g$.
    \textit{Under which conditions does the projective spectral entropy functional $\SPI[g]$
    produce an Einstein sector as its metric variation?}

    \medskip
    More precisely: define $A_g = -\nabla_g^2$ as the minimal admissible Laplace-type
    operator on a four-dimensional Riemannian manifold $(M,g)$ with a supplied metric, and
    \[
      \SPI[g] = \tfrac{1}{2}\log\det' A_g
    \]
    as the projective spectral entropy functional.
    The question is: what is the geometric content of $\delta\SPI/\delta g^{\mu\nu}$, in
    which regime does its local part reduce to the Einstein tensor, and which inputs must be
    supplied rather than derived?

    \medskip
    The answer has four stages, one per constituent paper, each with an explicit status.
    Gravity separates the regularization schemes, derives the decomposition of the
    renormalized variation given supplied renormalized coefficients, and states the Einstein-dominance and
    matching conditions.
    Thermodynamics derives the Einstein--matter coupling from spectral stationarity,
    conditionally on supplied coefficients and at leading local infrared order, and identifies the missing
    integrability bridge
    to any thermodynamic interpretation.
    Lorentz/CausalPropagation analyses the posited local four-derivative Lorentzian
    truncation and shows that its massless branch is exactly luminal.
    BornInfeld states the capacity axiom [A-cap] and presents the Born--Infeld structure as
    a saturation candidate; Gravity classifies the admissible tensorial completions and
    exhibits their non-uniqueness.

    \begin{remark}[Scope: gravity sector only]
This sub-programme concerns the gravitational dynamics in the absence of matter, plus the
conditional Einstein--matter coupling at the level of supplied functionals.
The gauge dynamics ($a_4$ Yang--Mills response from the vertical variation) belongs to
Note~8 (Gauge--Gravity Stratification Sub-Programme).
    \end{remark}

    \begin{remark}[Distinction from Note~2]
Note~2 (Emergent Geometry) addresses the \emph{reconstruction} of an effective metric from the admissibility
filter $\Pi_q$, with the target and the conditions stated above.
Note~4 (Spectral Gravity) addresses the \emph{dynamics} of the metric: the conditions under
which $\delta_g\SPI$ contains a dominant Einstein sector.
The two notes address logically distinct questions and should not be conflated.
    \end{remark}

  \section{Position in the Cosmochrony Programme}
    \label{sec:position}

    The spectral gravity sub-programme takes the following inputs.

    From Branch~I (axiomatic primitive): a genuinely emergent non-injective projection $\Pi$ has strictly positive
    projection entropy (ENI~\cite{Beau2026l}); the spectral functional $\tfrac{1}{2}\log\det' A_g$ used here is a
    sibling realisation motivated by the same information loss, and its identification with the projection
    entropy is open (ENI); the capacity axiom [A-cap] and the spectral length
    scale $\lsp$ are Branch~I objects (BornInfeld~\cite{Beau2026c}).

    On the geometric side: a four-dimensional Riemannian manifold $(M,g)$ with a supplied metric, on which the
    minimal Laplace-type operator $A_g = -\nabla_g^2$ is elliptic (Gravity~\cite{Beau2026h}).
    Its reading as the emergent metric of Note~2~\cite{Beau2026egnote} is conditional on [H-L], [H-lift] and a
    supplied rank-four extension; relating it to a Lorentzian metric needs a continuation that is not performed,
    and a Lorentzian counterpart of $\SPI$ requires an in-in construction that is open, as stated in the
    Lorentz/\allowbreak CausalPropagation paper~\cite{Beau2026i}.

    From Note~3 (gauge structure): the admissible principal bundle $P_{G_\Pi}(M, G_\Pi)$ for a supplied
    compact structure group, used in the gauge-extended functional $\SPI[g, A]$ of Q12~\cite{Beau2026q12} and Q13
    (Note~8 consumer).

    \medskip
    The outputs of the sub-programme feed directly into:
    \begin{itemize}
      \item Note~8 (Gauge--Gravity Stratification): the horizontal $a_2$ Einstein sector, under the conditions
      stated here, and the vertical $a_4$ Yang--Mills sector are obtained from the same functional, with all
      renormalized couplings as matching data.
      \item Q13 (Einstein--Yang--Mills synthesis): the conditional coupled system with
      matched coefficients.
      \item Cosmological and phenomenological applications, within the conditional scope
      stated in Section~\ref{sec:status}.
    \end{itemize}

  \section{Logical Chain of the Sub-Programme}
    \label{sec:chain}

    \begin{equation}
      \label{eq:chain}
      \underbrace{\SPI[g] = \tfrac{1}{2}\log\det' \Ag}_{\text{spectral entropy}}
      \;\Longrightarrow\;
      \underbrace{\delta_g\SPI^{\mathrm{ren}} \ni \cEH^{\mathrm{ren}}
        G_{\mu\nu}}_{\substack{a_2\text{ cutoff term} \\
        \text{IR dominant if } RL_4^2 \ll 1}}
      \;\Longrightarrow\;
      \underbrace{G_{\mu\nu} + \Lambda_{\mathrm{ren}} g_{\mu\nu} =
        8\pi\GN T_{\mu\nu}}_{\substack{\text{conditional stationarity} \\
        \text{of } \Gamma_\Pi = \SPI^{\mathrm{ren}} - W_\Pi}}
    \end{equation}
    \[
      \Longrightarrow\;
      \underbrace{\omega^2 = |\mathbf{k}|^2 \text{ (massless
        branch)}}_{\substack{\text{posited four-derivative} \\
        \text{truncation, exactly luminal}}}
      \;\Longrightarrow\;
      \underbrace{\sqrt{-\det(g + \lsp^2 R)} \in
        \{\Phi_a\}}_{\substack{\text{admissible BI member} \\
        \text{of a parametrised family}}}.
    \]

    The chain has four conceptually distinct stages.

    \begin{enumerate}
      \item \textit{Decomposition and conditional IR dominance.}
      The physical proper-time cutoff produces the power-sensitive local terms
      $a_0\Lambda^4$ and $a_2\Lambda^2$, while zeta regularization produces finite and
      logarithmic scale dependence governed by $a_4 = \zeta_A(0)$; the two schemes must not
      be conflated.
      The renormalized metric variation decomposes into Einstein, cosmological,
      higher-derivative, and non-local sectors, and the Einstein term dominates locally when
      the renormalized coefficient ratio satisfies $R L_4^2 \ll 1$.
      One minimal scalar contributes $\Delta\cEH = -\Lambda^2/[12(4\pi)^2]$: the finite
      renormalized value, its sign, and $\GN > 0$ require a matching condition and the
      complete operator content (Gravity~\cite{Beau2026h}).

      \item \textit{Conditional Einstein--matter coupling.}
      Given supplied renormalized coefficients and a projected matter functional
      $W_\Pi[g, \psi]$, stationarity of $\Gamma_\Pi = \SPI^{\mathrm{ren}} - W_\Pi$ yields
      $G_{\mu\nu} + \Lambda_{\mathrm{ren}} g_{\mu\nu} = 8\pi\GN T_{\mu\nu}$ at leading local infrared order.
      No local first law follows from the heat kernel alone: the spectral scale
      $u = -\partial_t\log K$ has the dimension of an eigenvalue, not of a temperature, and
      the integrability bridge to any thermodynamic reading is open
      (Thermodynamics~\cite{Beau2026Thermodynamics}).

      \item \textit{Lorentzian truncation.}
      On the posited local covariant four-derivative transverse-traceless kernel
      $\mathcal{O}_{\mathrm{TT}} = \Box(\cEH + \beta\Box)$, the kernel factorizes and the
      spectrum is the direct sum of an exactly luminal massless branch,
      $\omega^2 = |\mathbf{k}|^2$, and a massive spin-two pole with opposite residue.
      No quartic dispersion arises from this truncation, and no observational bound on
      $\lsp$ follows from it; a Lorentzian spectral action beyond the truncation requires an
      in-in construction, which is open (Lorentz/CausalPropagation~\cite{Beau2026i}).

      \item \textit{UV completion: admissibility without uniqueness.}
      Under conditions (C1)--(C5) and Hypothesis~[H-ext], the admissible tensorial
      completions form the family $\Phi(X) = \exp\sum_i h(\lambda_i)$ on the mixed Ricci carrier $X$, which the
      saturation condition applied to $X$ selects among linear two-derivative carriers, with $h$ constrained
      only at the origin and at the saturation boundary; the determinantal Born--Infeld
      density is the member $h(\lambda) = \tfrac12\log(1+\lambda)$, and the explicit family
      $\Phi_a(X) = \sqrt{\det(I+X)}\,e^{a\,\mathrm{tr}X^2}$ shows it is not uniquely
      selected (Gravity~\cite{Beau2026h}; BornInfeld~\cite{Beau2026c}).
    \end{enumerate}

  \section{Constituent Papers}
    \label{sec:papers}

    \subsection{Conditions for an infrared Einstein sector: Gravity}
      \label{sec:gravity}

      \textbf{Gravity}~\cite{Beau2026h} is the foundational paper of the sub-programme.
      It establishes the projective spectral entropy functional
      $\SPI[g] = \tfrac{1}{2}\log\det' A_g$ for a minimal Laplace-type operator on a
      four-dimensional Riemannian manifold and derives the renormalized metric variation via
      heat-kernel methods, separating the proper-time cutoff from zeta regularization.

      Its results are the following.

      \textit{Scheme separation and decomposition} (scheme separation proved; decomposition structural).
      The cutoff carries the power-sensitive terms $a_0\Lambda^4$ and $a_2\Lambda^2$; zeta
      regularization carries the finite and logarithmic terms governed by $a_4$.
      The renormalized variation decomposes as
      \[
        \frac{\delta\SPI^{\mathrm{ren}}}{\delta g^{\mu\nu}}
        = \cEH^{\mathrm{ren}}\, G_{\mu\nu}
        + c_\Lambda^{\mathrm{ren}}\, g_{\mu\nu}
        + \beta^{\mathrm{ren}}\, B_{\mu\nu}
        + T^{\mathrm{non\text{-}loc}}_{\mu\nu},
      \]
      and the Einstein term dominates locally when $R L_4^2 \ll 1$ with
      $L_4^2 = |\beta^{\mathrm{ren}}/\cEH^{\mathrm{ren}}|$.

      \textit{Matching status of the couplings} (structural).
      One minimal scalar contributes $\Delta\cEH = -\Lambda^2/[12(4\pi)^2]$, which is
      negative for the determinant convention used.
      The finite renormalized Einstein coefficient, its sign, and the observed Newton
      constant are matching data requiring the complete operator content and a
      renormalization condition; they are not predicted by the minimal determinant.

      \textit{Classification of the tensorial UV completions} (conditional on [H-ext]).
      Under conditions (C1)--(C5), the sector function $h$ is constrained only by
      $h(0) = 0$, $h'(0) = \tfrac12$, and divergence at the saturation boundary; every such
      $h$ defines an admissible completer.
      The determinantal Born--Infeld density is the member $h = \tfrac12\log(1+\lambda)$,
      and the one-parameter family $\Phi_a$ satisfies all conditions for every real $a$.
      Among linear two-derivative carriers, the infrared limit and the heat-kernel coefficient leave a trace
      parameter $\alpha$ of the spectral carrier free; the saturation condition, applied as stated to the curvature
      endomorphism, selects the mixed Ricci carrier $\alpha = 0$, and the non-uniqueness lies in $h$.
      The input a future uniqueness theorem would require (scalar-sector reduction,
      finite-order heat-kernel matching, or exceptional propagation) is stated explicitly.

    \subsection{Einstein--matter coupling from spectral stationarity: Thermodynamics}
      \label{sec:thermo}

      \textbf{Thermodynamics}~\cite{Beau2026Thermodynamics} derives the Einstein--matter
      coupling from a stationarity principle, and delimits what does not follow from the
      heat kernel.

      \textit{Dimensional status of the spectral scale} (proved).
      The diagonal heat-kernel scale $u(x;t) = -\partial_t\log K(x,x;t)$ has the dimension
      of a Laplacian eigenvalue ($L^{-2}$), not of a temperature ($L^{-1}$).
      No local first law and no multiplier field follow from the heat kernel alone; the
      integrability bridge --- a calibration, a heat one-form, and its integrability, in the
      sense of Carath\'eodory --- is an open problem.

      \textit{Conditional stationarity theorem} (conditional on supplied coefficients).
      Given the renormalized coefficients of Gravity and a projected matter functional
      $W_\Pi[g,\psi]$, stationarity of $\Gamma_\Pi = \SPI^{\mathrm{ren}} - W_\Pi$ yields, at leading local
      infrared order,
      \[
        G_{\mu\nu} + \Lambda_{\mathrm{ren}}\, g_{\mu\nu} = 8\pi\GN\, T_{\mu\nu},
      \]
      in Riemannian signature, on the Euclidean section.
      The derivation is variational and requires no horizon structure; its Lorentzian
      extension is open.

    \subsection{The four-derivative truncation: Lorentz/CausalPropagation}
      \label{sec:lorentz}

      \textbf{Lorentz/CausalPropagation}~\cite{Beau2026i} analyses the posited local
      covariant four-derivative transverse-traceless kernel of the linearized Lorentzian
      dynamics.

      \textit{Factorization and spectrum} (proved, given the posited kernel).
      For $\cEH \neq 0$, the kernel $\Box(\cEH + \beta\Box)$ factorizes and its kernel space
      is the direct sum $\ker\Box \oplus \ker(\cEH + \beta\Box)$.
      The massless branch is exactly luminal, $\omega^2 = |\mathbf{k}|^2$, with the standard
      two transverse-traceless helicity modes; the second branch is a massive spin-two pole
      with residue of opposite sign, tachyonic or not according to the sign of
      $-\cEH/\beta$.

      \textit{Consequences} (structural).
      No quartic graviton dispersion and no gravitational-wave bound on $\lsp$ follow from
      the truncation.
      The kernel itself is posited: a Lorentzian spectral action derived from $\SPI$
      requires an in-in (closed-time-path) construction, which is an open deliverable.

    \subsection{Capacity axiom and Born--Infeld candidate: BornInfeld}
      \label{sec:bi}

      \textbf{BornInfeld}~\cite{Beau2026c} states the capacity axiom [A-cap]: the per-node
      relaxation rate of an admissible relational configuration is uniformly bounded by
      $\cBI$.
      The axiom is motivated by non-injective projection together with finite local
      distinguishability; it is not a consequence of non-injectivity alone.
      Which continuum quantity inherits the bound is an open problem, and the Born--Infeld
      structure is presented as a saturation candidate, distinguished in the nonlinear
      electrodynamics literature by exceptional propagation rather than by saturation.
      Its conditional continuum statements (Lorentzian signature under a hyperbolicity
      hypothesis, the Newtonian exterior profile under spherical flux conservation) supply
      structural context for the present sub-programme; the parity structure of the
      candidate action feeds the O-series pair observable through O18, where the fibre
      identification is stated as an open problem.

      \begin{table}[ht]
        \centering
        \caption{Summary of constituent papers by stage.
        Status: P = proved (given stated inputs); S = structural; C = conditional; O = open.}
        \label{tab:papers}
        \renewcommand{\arraystretch}{1.2}
        \begin{tabular}{@{}llll@{}}
          \toprule
          Stage & Paper & Central output & Status \\
          \midrule
          Scheme separation & Gravity & cutoff ($a_0, a_2$) vs zeta ($a_4$) & P \\
          IR decomposition & Gravity & Einstein dominant if $\cEH \neq 0$, $RL_4^2 \ll 1$ & S \\
          Coupling status & Gravity & $\cEH^{\mathrm{ren}}$, $\GN$ are matching data & S \\
          UV completions & Gravity & admissible family; BI member not unique & C ([H-ext]) \\
          Spectral scale & Thermodynamics & $u = -\partial_t\log K$ is an eigenvalue scale & P \\
          Einstein--matter & Thermodynamics & conditional stationarity of $\Gamma_\Pi$ & C \\
          TT truncation & Lorentz/CausalProp & luminal massless branch; opposite-residue pole & P (posited kernel) \\
          Capacity axiom & BornInfeld & [A-cap]; BI as saturation candidate & Axiom / S \\
          \bottomrule
        \end{tabular}
      \end{table}

  \section{Inputs from Upstream Results}
    \label{sec:inputs}

    \begin{enumerate}
      \item \textbf{Non-injectivity as structural necessity (ENI~\cite{Beau2026l}).}
      Under the hypotheses of the ENI theorem, any genuinely emergent projection is non-injective and has strictly
      positive projection entropy.
      The spectral entropy functional is motivated by, not derived from, this result: ENI states that its
      identification with the projection entropy is open.

      \item \textbf{Supplied Riemannian metric $(M, g)$.}
      The sub-programme works on a four-dimensional Riemannian manifold with a supplied metric, on which the
      minimal Laplace-type operator $A_g$ is elliptic and its heat-kernel expansion is defined (Gravity).
      Its reading as the emergent metric of the emergent geometry sub-programme (Note~2~\cite{Beau2026egnote})
      is conditional.
      Under the unestablished [H-L] and the open [H-lift], Q5b~\cite{Beau2026q5b} gives a rank-three co-metric; a
      rank-four Lorentzian form requires a supplied extension, proportionality to $\eta^{\mu\nu}$ needs in addition
      the hypotheses of Q11~\cite{Beau2026q11} Corollary~6.1, relating the supplied metric to such a form needs a
      continuation from Riemannian signature that is not performed, and no metric coefficient value is established.

      \item \textbf{Spectral length scale $\lsp$ (Branch~I).}
      The scale $\lsp$ is the admissibility cutoff entering the power-sensitive
      contributions; it does not by itself fix the finite Newton constant, which is
      matching data.

      \item \textbf{Capacity axiom [A-cap] and candidate parity (BornInfeld, O18).}
      The capacity bound $\cBI$ is the axiom [A-cap] of~\cite{Beau2026c}; the parity
      structure of the Born--Infeld candidate action, established in the Weil setting by
      O18 as a response-family covariance, supplies the candidate parity used by the
      O-series pair observable.
      The identification of the parity orbit with the fibres of $\Pi$ is an open problem
      (O18), and the pair observable keeps the status of a structurally motivated
      hypothesis.
    \end{enumerate}

  \section{Outputs to Downstream Branches}
    \label{sec:outputs}

    \begin{table}[ht]
      \centering
      \caption{Principal outputs and downstream consumers.}
      \label{tab:outputs}
      \renewcommand{\arraystretch}{1.3}
      \begin{tabular}{@{}p{6.2cm}p{4.2cm}l@{}}
        \toprule
        Output & Consumer & Status \\
        \midrule
        Conditional Einstein sector (matched couplings) & Q13, Q6b & C \\
        $\cEH^{\mathrm{ren}}$, $\GN$ as matching data & Q12, Q13, Note~8 & S \\
        Admissible BI family (member not unique) & Q13 (conditional EYM) & C ([H-ext]) \\
        Luminal massless TT branch (posited kernel) & GW discussion & P \\
        Capacity axiom [A-cap]; candidate parity & O18, O30, O-series & Axiom / S \\
        $a_2 \to$ gravity; $a_4 \to$ gauge (stratification) & Note~8, Q12, Q13 & interpretive \\
        \bottomrule
      \end{tabular}
    \end{table}

  \section{Status of Results}
    \label{sec:status}

    \subsection{Proved results (given stated inputs)}
      \label{sec:proved}

      \begin{enumerate}
        \item \textit{Scheme separation} (Gravity):
        the proper-time cutoff carries $a_0\Lambda^4$ and $a_2\Lambda^2$; zeta
        regularization carries the finite and logarithmic $a_4$ terms; conflating the two
        schemes drops a factor $\Lambda^2$.

        \item \textit{One-scalar cutoff contribution} (Gravity):
        $\Delta\cEH = -\Lambda^2/[12(4\pi)^2]$ for the convention
        $\SPI = +\tfrac12\log\det' A_g$.

        \item \textit{Spectral carrier family} (Gravity; elementary, for carriers linear in the curvature, of
        two-derivative order and with a spectrum that is not a multiple of the identity):
        normalised by the Einstein limit, the carriers are
        $Y_\alpha = \tfrac{t}{d}\,I + X_0/(1-\alpha d)$, $\alpha \neq 1/d$, with $t = \mathrm{tr}\,X$ and $X_0$ the
        traceless part of $X = \lsp^2 g^{-1}R$; all give the same linear term, second-order data determine at most
        $|1-\alpha d|$, leaving the pair $\{\alpha,\, 2/d-\alpha\}$, and the carriers differ in their saturation
        boundary; within this classification $\alpha$ is free, and the saturation condition~(C3), applied as stated
        to the eigenvalues of $I + X$, selects $\alpha = 0$ among the completers $\exp\sum_i h(\mu_i)$ built on these
        carriers, with $\mu_i$ the eigenvalues of $Y_\alpha$ and $h$ a smooth sector function on $(-1,+\infty)$ with
        $h(0) = 0$, $h'(0) = \tfrac12$ and $h \to -\infty$ at $-1$ (Gravity, Theorem~1(o), which does not use
        [H-ext]).

        \item \textit{Constraints on the sector function} (Gravity; unconditional):
        conditions (C2)--(C3) fix only $h(0) = 0$, $h'(0) = \tfrac12$, and the boundary
        divergence; every such $h$ is admissible, whence non-uniqueness via the family
        $\Phi_a$.

        \item \textit{Dimensional status of the heat-kernel scale} (Thermodynamics):
        $u = -\partial_t\log K$ is an eigenvalue scale; it is not a temperature.

        \item \textit{Factorization of the posited TT kernel} (Lorentz/CausalPropagation):
        direct-sum spectrum; exactly luminal massless branch; opposite-residue massive
        spin-two pole.
      \end{enumerate}

    \subsection{Structural and conditional results}
      \label{sec:structural}

      \begin{enumerate}
        \item \textit{IR dominance of $G_{\mu\nu}$} (Gravity):
        the $a_2$ Einstein term dominates locally when the matched Einstein coefficient is non-zero and
        $RL_4^2 \ll 1$; a structural, power-counting statement about
        the renormalized decomposition, conditional on supplied coefficients, with no general bound on the
        non-local terms.

        \item \textit{Conditional Einstein--matter coupling} (Thermodynamics):
        stationarity of $\Gamma_\Pi$ yields
        $G_{\mu\nu} + \Lambda_{\mathrm{ren}} g_{\mu\nu} = 8\pi\GN T_{\mu\nu}$ on the
        Euclidean section at leading local infrared order, given matched coefficients and a supplied matter
        functional.

        \item \textit{Admissibility of the Born--Infeld completion} (Gravity; conditional
        on [H-ext]): the determinantal member belongs to the admissible family; its unique
        selection is an open problem with explicitly stated candidate inputs.

        \item \textit{Capacity axiom and candidate status} (BornInfeld):
        [A-cap] is an explicit axiom; the Born--Infeld structure is a saturation candidate;
        the constitutive, propagation, and gauge-sector bridges are open.
      \end{enumerate}

    \subsection{Open hypotheses and problems}
      \label{sec:open-hyp}

      \begin{enumerate}
        \item \textit{[H-ext] (admissible coherence extensivity)}.
        Its derivation from Axiom A3 (non-factorisability) is structurally motivated but
        not analytically closed; it conditions the classification of the tensorial
        completions.

        \item \textit{Matching principle for the couplings.}
        An independent principle selecting the positive finite Einstein coefficient ---
        rather than merely permitting it --- is the missing step before gravity can be
        claimed as an emergent spectral prediction.

        \item \textit{Integrability bridge to thermodynamics.}
        A calibration, a heat one-form, and its integrability would be required before any
        thermodynamic reading of $\SPI$; none is currently derived.

        \item \textit{In-in Lorentzian spectral action.}
        The four-derivative kernel is posited; deriving a Lorentzian spectral action from
        $\SPI$ requires a closed-time-path construction.

        \item \textit{Coupled field equations with matter in Lorentzian signature.}
        The construction of $T_{\mu\nu}^{(\Pi)}$ from projected matter configurations and
        the Lorentzian extension of the stationarity result remain open; the sourced
        vector sector is structurally fixed once a projected momentum density is supplied
        (frame dragging in the standard infrared regime), while its upstream derivation
        awaits the fermionic-matter sub-programme.

        \item \textit{Emergent reading of the metric.}
        The metric is supplied.
        Its reading as the emergent metric of Note~2 needs the spatial-limit hypothesis [H-L], the lifting
        hypothesis [H-lift], a rank-four extension of the rank-three co-metric of Q5b that no source constructs,
        a continuation from Riemannian signature that is not performed, and the data fixing the metric coefficients,
        which no source supplies.
      \end{enumerate}

  \section{Open Deliverables}
    \label{sec:open}

    Four open items define the boundary of the sub-programme.

    \textbf{Open deliverable 1: Matching principle for the renormalized couplings.}
    The sign and value of $\cEH^{\mathrm{ren}}$, hence $\GN$, must be selected by an
    independent principle involving the complete projected operator content.

    \textbf{Open deliverable 2: Analytical proof of [H-ext] and a selection principle for
    the completion.}
    [H-ext] conditions the classification; even under it, selecting the determinantal
    member requires input beyond conditions (C1)--(C5).

    \textbf{Open deliverable 3: In-in Lorentzian spectral action and Lorentzian
    stationarity.}
    A closed-time-path construction of the Lorentzian $\SPI$, and the extension of the
    conditional stationarity theorem beyond the Euclidean section.

    \textbf{Open deliverable 4: Coupled field equations with matter.}
    The full system $G_{\mu\nu} = 8\pi\GN T_{\mu\nu}$ with explicit matter fields requires
    identifying the admissible projected sector for each matter species, a task that awaits
    the completion of Note~6 (Fermionic Matter Sub-Programme).

  \section{Conclusion}
    \label{sec:conclusion}

    The spectral gravity sub-programme determines the conditions under which the projective
    spectral entropy functional $\SPI[g] = \tfrac{1}{2}\log\det' A_g$ supports gravitational
    dynamics.
    Throughout, the metric is supplied and Riemannian; its reading as the emergent metric of Note~2 is conditional
    on [H-L], [H-lift], a supplied rank-four extension and a continuation from Riemannian signature that is not
    performed, and no metric coefficient value
    is established.

    The renormalized metric variation contains an Einstein sector that dominates locally
    when $RL_4^2 \ll 1$, with all renormalized couplings --- including the sign of Newton's
    constant --- as matching data (Gravity).
    Conditionally on supplied coefficients and at leading local infrared order, spectral stationarity yields the
    Einstein--matter coupling on the Euclidean section, with the integrability bridge to any
    thermodynamic reading open (Thermodynamics).
    The posited four-derivative Lorentzian truncation propagates an exactly luminal massless
    graviton branch beside an opposite-residue massive pole, and yields no quartic
    dispersion and no observational bound on $\lsp$ (Lorentz/CausalPropagation).
    The determinantal Born--Infeld action is an admissible ultraviolet completion within an
    explicitly parametrised family, resting on the capacity axiom [A-cap] and awaiting a
    selection principle (Gravity, BornInfeld).

    The sub-programme does not address Yang--Mills dynamics or the gauge--gravity coupling
    ratio; those belong to Note~8.

    \medskip
    \textbf{Interpretive outlook.}
    The structural reading of the sub-programme is that spectral geometry fixes the menu of
    local gravitational structures and localises every remaining gap: a matching principle
    for the couplings, an integrability bridge for any thermodynamic reading, an in-in
    construction for the Lorentzian action, and a selection principle for the ultraviolet
    completion.
    Closing any one of these would convert the corresponding structural possibility into a
    derived result; until then, emergent spectral gravity is a precisely delimited
    programme, not an established conclusion.

    \newpage
    \bibliographystyle{plain}
    \bibliography{cosmochrony-bibliography}

\end{document}
